\ncertchapter{6}{Circular Kinematics}\label{ch:circkin}

Circular motion is where the vector nature of velocity stops being a formality and starts to matter. A
body can travel round a circle at perfectly constant speed and still be accelerating, and accelerating
hard, because its velocity, which is a vector, is continuously changing direction. That claim sounds
paradoxical until the definition of acceleration is taken seriously, and this chapter exists to take it
seriously.

The central result is that the acceleration of any body moving on a curve has exactly two parts. One,
along the velocity, changes the speed. The other, perpendicular to the velocity and pointing towards the
centre of the circle, changes the direction and has magnitude $v^2/r$, where $v$ is the speed and $r$ the radius of the circle. The chapter derives that result
carefully, from the product rule applied to a velocity written as speed times a unit vector, and then
explains why the second part is always perpendicular to the first. A great deal rests on it: every
problem in the dynamics of circular motion in \chref{ch:circdyn} begins from the radial acceleration found here.

Two further points are made along the way and are worth flagging. Angular velocity is a vector, but it
points along the axis of rotation rather than round the circle, and the reason is not arbitrary. And the
$v^2/r$ result is not special to circles: it applies at every point of \emph{any} curved path, with $r$
replaced by the local radius of curvature.

\section{The Problem Circular Motion Creates}

Everything so far has been handled by reducing motion to straight lines: even the projectile was split
into two fixed perpendicular directions, along each of which ``how fast'' told you nearly everything. On a
circle that strategy is awkward, because the direction of motion never stays fixed. A car going round a roundabout at a steady \SI{30}{km/h} has a
speedometer reading that never flickers, yet the passengers find themselves pressed against the doors,
the tyres are working hard, and something is happening that a straight-line description cannot capture.

\begin{motivbox}{The one sentence this chapter exists to unpack}
Velocity is a \textbf{vector}. A vector has two attributes, a \emph{size} and a \emph{direction},
and \emph{either} of them changing counts as a change in velocity. Acceleration is the rate of change of
velocity, so acceleration must be able to report on \textbf{both}. On a straight road only the size can
change, so one number suffices. On a circle the direction changes \emph{continuously and unavoidably},
and that second kind of change needs its own accounting. That accounting is the whole of this chapter.
\end{motivbox}

We build the vocabulary first (\S\ref{sec:angular}), then prove the central result
(\S\ref{sec:twoparts}), then classify what can happen, generalise it to any curved path, and treat
constant angular acceleration (\S\S\ref{sec:angleav}--\ref{sec:constalpha}).

\section{Angular Position, Velocity and Acceleration}\label{sec:angular}
A body moving on a circle is most naturally described by the angle it has turned through rather than by its
coordinates. This section defines the three angular quantities and connects each to its linear partner; it
begins by fixing the meaning of the three different lengths that appear in circular-motion problems.

\begin{remarkbox}{Notation used here and in the dynamics of circular motion}
Three different lengths appear in circular-motion problems and they must not be confused.
\begin{itemize}
\item $r$ is the radius of the \textbf{actual circular path} the particle travels. It is the length in
      the relations $v=r\omega$ and $v^{2}/r$ derived below, always measured perpendicularly from the \emph{axis of rotation}
      to the particle.
\item $R$ marks two special contexts, each stated where it arises: the \textbf{radius of curvature} of
      a general non-circular path in \S\ref{sec:curvature}, and the radius of the \textbf{vertical circle} in
      \chref{ch:circdyn}. In both it is still a path radius; the separate letter only flags the context.
      (In \chref{ch:projectiles} $R$ denoted the range; the context always makes clear which is meant.)
\item $\ell$ is the \textbf{length of a string or rod}, which is \emph{not} in general the radius of the
      circle. For a conical pendulum $r=\ell\sin\theta$. For a vertical circle the two happen to coincide,
      which is why that case is written with one symbol.
\end{itemize}
For angles, $\theta$ is the general position or inclination angle, $\phi$ the angle the total
acceleration makes with the tangent, and, in later chapters, $\lambda$ the angle of friction and
$\beta$ the angle at which a force is applied. The symbol $\alpha$ is reserved \emph{throughout} for \textbf{angular acceleration}
and is never used for an angle.
\end{remarkbox}

A particle on a circle of fixed radius has only \emph{one} degree of freedom: once you know the angle,
you know where it is. So it is wasteful to track $x$ and $y$ separately; one angle does the job.

\begin{defbox}{The three angular quantities}
\begin{itemize}
\item \textbf{Angular position} $\theta$: the angle between the radius to the particle and a chosen
      reference direction, measured in \textbf{radians} and, by convention, positive anticlockwise.
\item \textbf{Angular velocity} $\omega=\dfrac{d\theta}{dt}$ (\si{rad/s}): how fast the radius sweeps,
      positive for anticlockwise turning. It is treated here as a signed scalar; its vector form
      comes later in this section.
\item \textbf{Angular acceleration} $\alpha=\dfrac{d\omega}{dt}=\dfrac{d^{2}\theta}{dt^{2}}$
      (\si{rad/s^2}): how fast the sweeping rate itself changes.
\end{itemize}
\end{defbox}

\why{\textbf{Why radians, and not degrees?} Because the radian is \emph{defined} by $\theta=s/r$, with $s$ the length of arc swept out (\chref{ch:toolkit}), so it
converts an angle into a length the instant you multiply by $r$: $s=r\theta$. Degrees carry an arbitrary
factor of $\pi/180$ that would then appear in every formula. Above all,
$\dfrac{d}{d\theta}\sin\theta=\cos\theta$ holds only when $\theta$ is in radians (\chref{ch:toolkit}), so every formula below that contains $\theta$, $\omega$ or $\alpha$ requires
radians. Degrees remain fine for \emph{naming} an angle, as in $\sin90^\circ$.}

\begin{figure}[htbp]
\centering
\tikzfig{cm_angular}
\figcaption{fig:cmang}{The three angular quantities and their linear partners. Each bridge is the one above it, differentiated once: $s=r\theta$ gives $v=r\omega$, which gives $a_t=r\alpha$.}
\end{figure}

\begin{derivbox}{The three bridges to linear quantities}
Each bridge is one differentiation of the one above it; Figure~\ref{fig:cmang} marks each angular
quantity on the circle beside its linear partner.

\stepn{1} \textbf{Position.} By the definition of the radian, arc $=$ radius $\times$ angle:
\[ s=r\theta. \]

\stepn{2} \textbf{Velocity.} Differentiate with $r$ held constant (the particle stays on the circle):
\[
\frac{ds}{dt}=r\frac{d\theta}{dt}
\qquad\Longrightarrow\qquad
\boxed{\;v=r\omega\;}
\]

\stepn{3} \textbf{Acceleration (the tangential part).} Differentiate once more:
\[
\frac{dv}{dt}=r\frac{d\omega}{dt}
\qquad\Longrightarrow\qquad
\boxed{\;a_t=r\alpha\;}
\]
\textbf{Read this carefully:} $a_t$ is the rate of change of \emph{speed}, the rate at which your
speedometer reading changes. It is \emph{not} the whole acceleration, and \S\ref{sec:twoparts} explains what is missing.
\end{derivbox}

Every point of a rigid wheel turns through the same angle in the same time, so all of them share one
$\omega$; but $v=r\omega$ grows with distance from the axle. The rim of a spinning wheel moves faster
than the hub, while every spoke sweeps at the same angular rate. That is why $\omega$, not $v$, is the
natural label for a rotation as a whole.

When the motion repeats at a steady rate, it is often described by how long one turn takes or how many
turns occur per second. One revolution is $2\pi$ radians, which gives the following relations.

\begin{keybox}{Period and frequency: the small relations that get forgotten}
For uniform circular motion, with $T$ the \textbf{period} (time for one revolution, in seconds; not
the time of flight that $T$ denoted in \chref{ch:projectiles}),
\[
T=\frac{2\pi}{\omega}=\frac{1}{\nu},
\qquad
\omega=2\pi\nu,
\qquad
T=\frac{2\pi r}{v},
\]
where $\nu$ is the frequency in hertz (revolutions per second); the middle relation follows because one
revolution, $2\pi$ radians, takes time $T$. For a rate quoted as $N$ revolutions per minute (\textbf{rpm}),
divide by $60$ first: $\nu=\dfrac{N}{60}$ and $\omega=\dfrac{2\pi N}{60}$.
\end{keybox}

\subsection{Angular velocity is a vector, and it points the strange way}

Here is a puzzle. Angular velocity has a magnitude and a sense (clockwise or anticlockwise), so it
ought to be a vector. But which way does the arrow point? Not along the motion: the particle's
direction changes every instant, so any arrow following it would be useless as a label for the
\emph{whole} rotation.

For rotation about a fixed axis, the only direction that is \emph{the same for every particle of the
rotating body, and constant in time}, is the one perpendicular to the plane of rotation: the
\textbf{axis}. So that is where we point the arrow. The sense along the axis is fixed by the right-hand
rule of Figure~\ref{fig:omegavec}, a convention chosen so that $\vv=\vomega\ex\vr$ comes out right, where $\vomega$ is the angular velocity vector (magnitude $\omega$, in \si{rad/s}) and $\vr$ is the position vector of the particle (in metres) measured from a point on the axis. The cross product is the one of \chref{ch:vectors}.


\begin{figure}[htbp]
\centering
\tikzfig{cm_omegavec}
\figcaption{fig:omegavec}{Curl the fingers of the right hand in the sense of rotation; the thumb gives
$\vomega$. Drawn on the page, it is either into the page ($\otimes$, clockwise) or out of it ($\odot$,
anticlockwise). Quantities lying \emph{in} the plane are $\vv$, $\va$ and the two components of $\va$ found in \S\ref{sec:twoparts};
those perpendicular to it are $\vomega$ and $\vc{\alpha}$.}
\end{figure}

\begin{keybox}{The vector relation}
\[
\boxed{\;\vv=\vomega\ex\vr\;}
\]
Check that it works, with the origin at the centre of the circle:
$|\vv|=\omega r\sin90^\circ=\omega r$ \checkmark, and the cross product is
automatically perpendicular to $\vr$; that is, tangential \checkmark. The vector form also gives the
\emph{sense} of the motion, which the scalar $v=\omega r$ cannot.

\medskip
$\vc{\alpha}$ likewise points along the axis. If the body is \textbf{speeding up}, $\vc{\alpha}$ is
\emph{parallel} to $\vomega$; if \textbf{slowing down}, $\vc{\alpha}$ is \emph{antiparallel}. There is no
other option: for rotation about a fixed axis both must lie along that axis.
\end{keybox}

\section{Why the Acceleration Must Have Exactly Two Parts}\label{sec:twoparts}

The bridge $a_t=r\alpha$ accounts only for changes in \emph{speed}. Yet the car on the roundabout
has a steady speed and its velocity still changes, because its direction does. Something is missing
from the acceleration, and this section finds it. This is the heart of the chapter; what follows it
applies the result.

Write the velocity as \emph{size times direction}:
\[
\vv \;=\; \underbrace{v}_{\text{how fast}}\;\times\;\underbrace{\uhat{t}}_{\text{which way}}
\]
where $\uhat{t}$ is a unit vector pointing along the motion. A vector in a plane is completely
specified by those two numbers, a size and a direction (an angle), and by nothing else. So when we ask
how the velocity can change, there are exactly two answers available.

\smallskip
\textbf{A familiar version of the same statement.} A driver has two kinds of control that alter the
car's velocity. The accelerator and brake change \emph{how fast} you are going. The steering
wheel changes \emph{which way} you are going. They are independent: you may use either alone, or both
together, and between them they exhaust what you can do to the car's velocity. The two components of
acceleration are nothing more than the physicist's names for those two controls.


\begin{motivbox}{Why the two changes are perpendicular: watch the tip of the arrow}
Draw the velocity as an arrow and follow its \emph{tip} over a short interval (Figure~\ref{fig:twoways}).

\smallskip
\textbf{If only the size changes}, the arrow grows or shrinks but keeps its line, so the tip slides
\emph{along} the arrow's own direction.

\smallskip
\textbf{If only the direction changes}, the length is fixed, so the tip is confined to a circle of
radius $v$ centred on the arrow's tail. A point moving on a circle moves along the tangent to that
circle, and the tangent to a circle is perpendicular to its radius. The radius here \emph{is} the
arrow. So the tip moves \emph{across} the arrow.

\smallskip
One motion along, one across. They are perpendicular not by assumption or by convenient choice, but
because that is the only way the tip of a vector can move. Add them and you have the general case, in
which the acceleration has a component that changes the speed and a component that bends the path.
\end{motivbox}

\begin{figure}[htbp]
\centering
\tikzfig{cm_twoways}
\figcaption{fig:twoways}{The whole argument in one picture. A change in size moves the tip along the
arrow; a change in direction moves it across; a general change does both. Because the second motion is
confined to a circle of fixed radius, it is always perpendicular to the first.}
\end{figure}

The picture shows \emph{that} there are two perpendicular parts; it does not say how big the sideways
part is. The product rule settles both questions at once.

\begin{derivbox}{The two components, from the product rule}
\stepn{1} \textbf{Differentiate the product.} With $\vv=v\,\uhat{t}$,
\[
\va=\frac{d\vv}{dt}=\underbrace{\frac{dv}{dt}\,\uhat{t}}_{\text{size changing}}
\;+\;\underbrace{v\,\frac{d\uhat{t}}{dt}}_{\text{direction changing}}.
\]
Two terms, because a product of two things can change in two ways. Nothing has been assumed about
circles yet; this holds for \emph{any} path.

\stepn{2} \textbf{Prove the second term is perpendicular to the first.} Because $\uhat{t}$ is a unit
vector, $\uhat{t}\cdot\uhat{t}=1$ at all times. Differentiate both sides:
\[
2\,\uhat{t}\cdot\frac{d\uhat{t}}{dt}=0
\qquad\Longrightarrow\qquad
\frac{d\uhat{t}}{dt}\perp\uhat{t}.
\]
This short result carries much of the argument: it says the change of direction \emph{must} be
sideways. It is the formal version of ``a fixed-length arrow can only swing''.

\stepn{3} \textbf{Evaluate how fast the direction turns.} Two things must be established here, and
neither is obvious: \emph{through what angle} $\uhat{t}$ turns, and \emph{how far its tip moves} as a
result.

\smallskip
\textbf{(3a) The tangent turns through the same angle as the radius.} The tangent to a circle is
perpendicular to the radius at the point of contact, and it stays perpendicular as the particle moves.
So the radius and the tangent form a rigid right angle that is carried around the circle together, as
in Figure~\ref{fig:dthat}(a). If
the radius swings through $d\theta$, a line locked at $90^\circ$ to it has no choice but to swing
through $d\theta$ as well. Hence $\uhat{t}$ turns through exactly $d\theta$, the same angle that
appears in $\omega=d\theta/dt$.

\smallskip
\textbf{(3b) How far the tip of a unit vector moves when it turns.} Draw $\uhat{t}(t)$ and
$\uhat{t}(t+dt)$ \emph{from a common tail}, as in Figure~\ref{fig:dthat}(b). Both have length exactly $1$, because they are unit
vectors, and the angle between them is $d\theta$ by part (3a). They therefore form an
\textbf{isosceles triangle} with two unit sides and apex angle $d\theta$, and the third side of that
triangle is precisely the change $d\uhat{t}$. Dropping a perpendicular from the apex bisects both the
angle and the base, giving
\[
\left|d\uhat{t}\right| = 2\,(1)\sin\frac{d\theta}{2}.
\]
For an infinitesimal angle $\sin(d\theta/2)\to d\theta/2$ (the small-angle approximation of
\chref{ch:toolkit}), so
\[
\left|d\uhat{t}\right| = d\theta.
\]
\emph{An equivalent one-line argument:} the tip of a unit vector is confined to a circle of radius $1$,
and when the vector turns by $d\theta$ the tip travels an arc of length $1\times d\theta$; for a small
angle the chord equals the arc.

\smallskip
\textbf{(3c) Therefore the rate.} Dividing by $dt$,
\[
\left|\frac{d\uhat{t}}{dt}\right|=\frac{d\theta}{dt}=\omega.
\]

\smallskip
\textbf{(3d) And its direction.} In the isosceles triangle the base is perpendicular to the bisector of
the apex angle. As $d\theta\to0$ that bisector becomes $\uhat{t}$ itself, so $d\uhat{t}$ ends up
perpendicular to $\uhat{t}$, confirming Step 2 geometrically. Of the two perpendicular directions it
leans towards the \emph{inside} of the turn, because that is the side the vector is swinging towards.
Writing $\uhat{n}$ for the inward (centre-pointing) unit vector,
\[
\frac{d\uhat{t}}{dt}=\omega\,\uhat{n}.
\]
Here, and wherever $\omega$ multiplies the speed below, $\omega$ means the angular \emph{speed}
$|d\theta/dt|$, so the coefficient is never negative.

\Needspace*{27\baselineskip}
\begin{figbox}
\centering
\tikzfig{cm_dthat}
\end{figbox}
\figcaption{fig:dthat}{The two halves of Step 3. \textbf{(a)} The tangent is locked at $90^\circ$ to the
radius, so it turns through the same $d\theta$. \textbf{(b)} Placing the two unit tangents tail to tail
gives an isosceles triangle whose equal sides are $1$ and whose apex angle is $d\theta$; its base is
$d\uhat{t}$, of length $2\sin(d\theta/2)\approx d\theta$.}

\stepn{4} \textbf{Assemble.}
\[
\boxed{\;\va=\underbrace{\frac{dv}{dt}}_{a_t}\uhat{t}\;+\;\underbrace{v\omega}_{a_r}\uhat{n}\;}
\]
and using $\omega=v/r$ the second coefficient becomes the familiar
\[
a_r=v\omega=\frac{v^{2}}{r}=\omega^{2}r.
\]
\end{derivbox}

\begin{tablebox}{What each component is \emph{for}}
\begin{center}
\begin{tabular}{@{}lll@{}}
\toprule
& \textbf{Tangential} $a_t$ & \textbf{Radial (centripetal)} $a_r$\\
\midrule
Changes & the \emph{speed} & the \emph{direction}\\
Direction & along $\vv$ (or against it) & always towards the centre\\
Formula & $\dfrac{dv}{dt}=r\alpha$ & $\dfrac{v^{2}}{r}=\omega^{2}r$\\
Can it be zero? & \textbf{Yes}, whenever the speed is steady & \textbf{Only} at an instant when $v=0$\\
\bottomrule
\end{tabular}
\end{center}
\end{tablebox}

\begin{trapbox}{$\alpha$ does \emph{not} have two components}
It is easy to slide from ``the acceleration has two components'' into ``the angular acceleration has
two components''. It does not, and keeping the two words apart matters.
\begin{itemize}
\item \textbf{Linear acceleration $\va$} is a vector lying \emph{in} the plane of the circle. It has
      two components, $a_t$ and $a_r$, because a planar vector has two degrees of freedom.
\item \textbf{Angular acceleration $\vc{\alpha}$} is a single axial quantity, pointing \emph{along the
      rotation axis}, parallel or antiparallel to $\vomega$. For rotation about a fixed axis there is
      no second direction available to it, so it has nothing to split into.
\end{itemize}
The relationship between them is one-way: $\alpha$ feeds the tangential component only, through
$a_t=r\alpha$. The radial component $a_r=v^{2}/r$ is present even when $\alpha=0$, which is precisely
the case of uniform circular motion: no angular acceleration at all, yet a perfectly real linear
acceleration pointing at the centre.
\end{trapbox}

\why{\textbf{Why a moving body on a curve cannot have $a_r=0$.} Since $a_r=v\,|d\uhat{t}/dt|$, while the
body is moving ($v\neq0$) a zero $a_r$ would require $d\uhat{t}/dt=\vc0$: the direction of travel would be
frozen, and the body would be moving in a \emph{straight line}. A body on a circle, by definition, is not. So the moment a path curves at all, there
\emph{must} be an inward acceleration, and the sharper the curve (smaller $r$) or the faster the
travel (larger $v$), the larger it must be. This is why $a_r=v^{2}/r$ grows with $v^{2}$ and with $1/r$:
you are asking the velocity vector to swing round faster on both counts.

Turned round, this becomes the most useful sentence in the dynamics of circular motion
(\chref{ch:circdyn}): \emph{curving is not something that happens to a body; it is something that must
be continuously caused.}}

\begin{exbox}{The car on the roundabout}
The car from the start of the chapter, at a steady \SI{30}{km/h}, goes round a roundabout of
radius \SI{20}{m}. Find its acceleration.

\stepn{1} \textbf{Convert to SI.} $v=30/3.6\approx\SI{8.33}{m/s}$.

\stepn{2} \textbf{Radial part.} The speed is steady, so $a_t=0$ and the whole acceleration is
$a_r=v^2/r=69.4/20\approx\SI{3.5}{m/s^2}$, directed towards the centre of the roundabout.

That is about a third of $g$, with the speedometer never moving. Doubling the speed would quadruple it.
\ans{$a\approx\SI{3.5}{m/s^2}$, towards the centre}
\end{exbox}

\begin{trapbox}{there is no outward acceleration}
``Centripetal'' means \emph{centre-seeking}; it names the direction of $a_r$, not a new kind of
acceleration. The passengers on the roundabout feel pressed \emph{outwards} against the door, but
nothing is accelerating them outwards. Left to themselves they would carry on in a straight line
(\chref{ch:newton}); the car turns underneath them, and the door, pushing \emph{inwards}, is what makes
them follow the curve. Measured from the ground, the acceleration of every part of the car and its
contents points towards the centre. The outward ``centrifugal'' effect appears only when the motion is
described from inside the turning car, a description \chref{ch:circdyn} handles with care.
\end{trapbox}

\begin{keybox}{Total acceleration}
The two components are perpendicular, so they combine by Pythagoras:
\[
\boxed{\;a=\sqrt{a_t^{2}+a_r^{2}}=\sqrt{\left(\frac{dv}{dt}\right)^{2}+\left(\frac{v^{2}}{r}\right)^{2}}\;}
\qquad
\tan\phi=\frac{a_r}{a_t}
\]
where $\phi$ is measured from the direction of motion, so that $\phi$ is the angle between $\va$ and
$\vv$ ($0\le\phi\le180^\circ$), with $a_t=a\cos\phi$ and $a_r=a\sin\phi$.
\end{keybox}

\section{\texorpdfstring{Reading the Angle Between $\va$ and $\vv$}{Reading the Angle Between a and v}}\label{sec:angleav}

The previous section split the acceleration into two perpendicular parts and gave their sizes. The
angle between the total acceleration and the velocity therefore encodes how the two parts compare, and
reading it is a quick way to classify a motion. You can diagnose any circular motion from a single picture, without a formula, by looking at where the
acceleration arrow points relative to the velocity arrow, as in the three panels of
Figure~\ref{fig:threetypes}. Since $a_t=a\cos\phi$:
\begin{itemize}
\item $\phi<90^\circ$ (acute): $\cos\phi>0$, so $a_t>0$: the acceleration has a forward part.
      \textbf{Speeding up.}
\item $\phi=90^\circ$: $\cos\phi=0$, so $a_t=0$: no forward or backward part at all.
      \textbf{Constant speed.}
\item $\phi>90^\circ$ (obtuse): $\cos\phi<0$, so $a_t<0$: the acceleration has a backward part.
      \textbf{Slowing down.}
\end{itemize}
Notice the middle case: the acceleration is entirely perpendicular to the motion, which is exactly the
condition for the speed to be unchanged. \chref{ch:energy} will meet the same fact in another form: a
force perpendicular to the motion does no work.


\begin{figure}[htbp]
\centering
\tikzfig{cm_threetypes}
\figcaption{fig:threetypes}{The three kinds of circular motion, told apart by the angle between $\va$ and
$\vv$. In every case the radial component is present and points to the centre; what changes is whether a
tangential component accompanies it, and which way it faces.}
\end{figure}

\begin{keybox}{The three types, in full}
\begin{enumerate}
\item \textbf{Uniform.} $v$ and $\omega$ constant $\Rightarrow a_t=0$, $\alpha=0$. The acceleration is
      purely radial, $a=a_r=v^{2}/r$, perpendicular to $\vv$ at all times.
\item \textbf{Accelerated.} $v$ and $\omega$ increasing. $a_t$ points along $\vv$; $\vc{\alpha}$ is
      \emph{parallel} to $\vomega$; the angle between $\va$ and $\vv$ is acute.
\item \textbf{Retarded.} $v$ and $\omega$ decreasing. $a_t$ points against $\vv$; $\vc{\alpha}$ is
      \emph{antiparallel} to $\vomega$; the angle is obtuse.
\end{enumerate}
In cases 2 and 3 the motion is \emph{non-uniform} circular motion, and $a=\sqrt{a_t^{2}+a_r^{2}}$ must
be used; quoting $v^{2}/r$ alone is then an incomplete answer.
\end{keybox}

\section{Any Curved Path: The Radius of Curvature}\label{sec:curvature}

Look back at the derivation: Steps 1 and 2 used only that $\vv=v\uhat{t}$ with $\uhat{t}$ a unit vector.
That is true for \emph{any} motion whatsoever: a projectile, a roller-coaster, a swerving car. So
every curved path has a tangential acceleration and a component perpendicular to the motion. On a
general path there is no fixed centre, so this second component is called the \textbf{normal}
acceleration $a_n$ rather than the radial one. The formula $a_n=v^{2}/R$ survives intact, provided we
ask: \emph{what is $R$ for a path that is not a circle?}

The answer is the natural one. At each point, a curve is locally indistinguishable from \emph{some}
circle: the one that hugs it most closely there, drawn in Figure~\ref{fig:curvi} with its centre marked. That circle's radius is the \textbf{radius of
curvature} $R$ at that point. A gentle bend has large $R$; a hairpin has small $R$; a straight line has
$R=\infty$, which correctly gives $a_n=0$.

To see that $a_n=v^{2}/R$ really does hold, repeat Step~3 of the derivation on that best-fitting circle.
Let $\psi$ be the angle the tangent makes with a fixed direction. As the body moves a distance $ds$ along
the path, the tangent turns through $d\psi=ds/R$; this relation is in fact the precise definition of
$R$, namely $R=ds/d\psi$. Then $|d\uhat{t}/dt|=d\psi/dt=(ds/dt)/R=v/R$, and the second term of Step~1 has
magnitude $v\cdot v/R=v^{2}/R$.


\begin{figure}[htbp]
\centering
\tikzfig{cm_curvilinear}
\figcaption{fig:curvi}{On any curved path, resolve $\va$ along and across $\vv$. The along-part changes
the speed; the across-part bends the path and equals $v^{2}/R$, where $R$ is the radius of the circle
that best fits the curve at that point. Its centre lies on the concave side.}
\end{figure}

\begin{methodbox}{Working backwards to find $R$}
Given the acceleration and velocity at an instant, the radius of curvature follows without any geometry.
The dot product picks out the part of $\va$ along $\vv$; Pythagoras gives the rest; and $a_n=v^2/R$
is turned round:
\[
a_t=\frac{\va\cdot\vv}{v},\qquad
a_n=\sqrt{a^{2}-a_t^{2}},\qquad
\boxed{\;R=\frac{v^{2}}{a_n}\;}
\]
\textbf{Worked instance: a projectile at its apex.} There $\vv$ is horizontal, with $v=v_0\cos\theta$, and
the acceleration is the free-fall acceleration $g$ of \chref{ch:projectiles}, vertically downward, so they are perpendicular: $a_t=0$ and $a_n=g$. Hence
\[
R=\frac{v^{2}}{g}=\frac{v_0^{2}\cos^{2}\theta}{g},
\]
the radius of the circle that best fits the parabola at the top of its flight. A projectile is
\emph{not} in circular motion, yet the method applies without change.
\end{methodbox}

\section{Equations for Constant Angular Acceleration}\label{sec:constalpha}

So far $\alpha$ has appeared only as the source of $a_t$. Many problems instead ask how a wheel or fan
gets from one angular velocity to another: how long it takes, and how many turns it makes on the way.
That needs $\omega$ and $\theta$ as functions of time. When the angular acceleration is constant, the angular quantities obey equations of exactly the same form as
the kinematic equations of \chref{ch:straightline}, and under the same strict condition.

\begin{keybox}{The constant-$\alpha$ equations, and when they are legal}
When $\alpha$ is \textbf{constant}, and only then, the angular equations are exact copies of the linear ones
(the derivations are identical, with $\theta,\omega,\alpha$ in place of $s,v,a$):
\[
\omega=\omega_0+\alpha t,\qquad
\Delta\theta=\omega_0t+\tfrac12\alpha t^{2},\qquad
\omega^{2}=\omega_0^{2}+2\alpha\,\Delta\theta,
\]
where $\omega_0$ and $\theta_0$ are the values at $t=0$ and $\Delta\theta=\theta-\theta_0$ is the angle
turned (in radians).
\end{keybox}

\begin{methodbox}{When $\alpha$ is not constant, go back to the definitions}
The formulas above are then invalid, and you integrate instead:
\[
\omega=\frac{d\theta}{dt},\qquad
\alpha=\frac{d\omega}{dt},\qquad
\text{so}\qquad
\int d\theta=\int\omega\,dt,\qquad \int d\omega=\int\alpha\,dt.
\]
If $\alpha$ is given as a function of \emph{angle} rather than time, use the chain-rule form
\[
\boxed{\;\alpha=\frac{d\omega}{dt}=\frac{d\omega}{d\theta}\cdot\frac{d\theta}{dt}=\omega\frac{d\omega}{d\theta}\;}
\qquad\Longrightarrow\qquad
\int\omega\,d\omega=\int\alpha\,d\theta,
\]
the exact rotational twin of the $v\,dv/dx$ trick from \chref{ch:straightline}.
\end{methodbox}

\begin{exbox}{An angular acceleration that grows with time}
A wheel starts from rest with $\alpha=6t$ (\si{rad/s^2}, $t$ in seconds). Find $\omega$ and the angle
turned at $t=\SI{2}{s}$.

\stepn{1} \textbf{Integrate $\alpha$.} $\omega=\int_0^t 6t\,dt=3t^2$, so $\omega=\SI{12}{rad/s}$ at
$t=\SI{2}{s}$.

\stepn{2} \textbf{Integrate $\omega$.} $\Delta\theta=\int_0^t 3t^2\,dt=t^3=\SI{8}{rad}$, about $1.3$
revolutions.

The constant-$\alpha$ formula $\tfrac12\alpha t^2$, with the final $\alpha=\SI{12}{rad/s^2}$, would give
\SI{24}{rad}, three times too large.
\ans{$\omega=\SI{12}{rad/s}$; $\Delta\theta=\SI{8}{rad}$}
\end{exbox}

\begin{exbox}{A speeding-up particle}
A particle moves on a circle of radius \SI{0.5}{m}, its speed rising uniformly from \SI{2.0}{m/s} to
\SI{4.0}{m/s} in \SI{4.0}{s}. Find the angular acceleration, and the total acceleration at the moment the
speed is \SI{4.0}{m/s}.

\stepn{1} \textbf{Tangential acceleration}: the rate of change of speed:
\[
a_t=\frac{\Delta v}{\Delta t}=\frac{4.0-2.0}{4.0}=\SI{0.5}{m/s^2}.
\]

\stepn{2} \textbf{Angular acceleration} from $a_t=r\alpha$:
\[
\alpha=\frac{a_t}{r}=\frac{0.5}{0.5}=\SI{1}{rad/s^2}.
\]

\stepn{3} \textbf{Radial acceleration} at $v=\SI{4.0}{m/s}$:
\[
a_r=\frac{v^{2}}{r}=\frac{16}{0.5}=\SI{32}{m/s^2}.
\]

\stepn{4} \textbf{Total.} The two are perpendicular:
\[
a=\sqrt{a_t^{2}+a_r^{2}}=\sqrt{0.25+1024}\approx\SI{32.0}{m/s^2}.
\]
\textbf{Notice how lopsided this is:} the radial part is $64$ times the tangential part. On a tight
curve, \emph{bending the path can cost far more acceleration than changing the speed}; quoting $a_t$
alone would miss almost all of it.
\ans{$\alpha=\SI{1}{rad/s^2}$; $a\approx\SI{32}{m/s^2}$}
\end{exbox}

\begin{exbox}{When the speed is given as a function of time}
A particle moves on a circle of radius \SI{2}{m} with $v=t^{2}$ (SI units). Find the radial, tangential
and net accelerations at $t=\SI{2}{s}$.

\stepn{1} \textbf{Speed at that instant:} $v=(2)^{2}=\SI{4}{m/s}$.

\stepn{2} \textbf{Radial:} $a_r=\dfrac{v^{2}}{r}=\dfrac{16}{2}=\SI{8}{m/s^2}$.

\stepn{3} \textbf{Tangential:} differentiate the \emph{speed}, not the position:
$a_t=\dfrac{dv}{dt}=2t=\SI{4}{m/s^2}$ at $t=2$.

\stepn{4} \textbf{Net:} $a=\sqrt{8^{2}+4^{2}}=\sqrt{80}\approx\SI{8.94}{m/s^2}$, making an angle
$\tan^{-1}(8/4)\approx63.4^\circ$ with the tangent: an \emph{acute} angle with $\vv$, confirming that
the particle is speeding up. \ans{$a_r=\SI{8}{m/s^2}$, $a_t=\SI{4}{m/s^2}$, $a=\sqrt{80}\approx\SI{8.94}{m/s^2}$}
\end{exbox}

\begin{trapbox}{$a_t$ is the derivative of the \emph{speed}, not of the velocity}
$\dfrac{d|\vv|}{dt}$ and $\left|\dfrac{d\vv}{dt}\right|$ are different quantities. The first is $a_t$
alone; the second is the \emph{whole} acceleration $a$. In uniform circular motion the first is zero
while the second is $v^{2}/r$. Confusing the two is the usual source of the belief that uniform
circular motion is unaccelerated.
\end{trapbox}

\medskip
\subsection{A conceptual test: which products can be negative?}

Nothing here needs a calculator, yet it checks whether you have absorbed \emph{which vectors lie in the
plane of the circle and which stand perpendicular to it}; Figure~\ref{fig:signs} draws the vectors in
each product. Students who can compute $v^{2}/r$ fluently
often cannot answer it, a sign that the geometry has not yet been understood.


\begin{figure}[htbp]
\centering
\tikzfig{cm_signs}
\figcaption{fig:signs}{The three dot products. The answer in each case follows from one fact: $\vv$,
$\va$, $a_t$ and $a_r$ lie \emph{in} the plane of the circle, while $\vomega$ and $\vc{\alpha}$ stand
\emph{perpendicular} to it, along the axis.}
\end{figure}

\begin{exbox}{Signs of $\vomega\cdot\vv$, $\vv\cdot\va$ and $\vomega\cdot\vc{\alpha}$}
For a particle in circular motion, state whether each product can be positive, negative or zero.

\stepn{1} \textbf{$\vomega\cdot\vv$, always zero.} $\vomega$ points along the axis, perpendicular to
the plane of the circle; $\vv$ is tangential and therefore lies \emph{in} that plane. Two perpendicular
vectors have zero dot product, and this holds at every instant, in every kind of circular motion.
\[
\boxed{\vomega\cdot\vv=0\ \text{always}}
\]

\stepn{2} \textbf{$\vv\cdot\va$, can be any of the three.} Both lie in the plane, so the angle
between them is unconstrained. From \S\ref{sec:angleav}:
\begin{itemize}
\item \emph{positive} when the angle is acute: the particle is speeding up;
\item \emph{zero} when it is $90^\circ$ (uniform circular motion), or at an instant when $v=0$;
\item \emph{negative} when it is obtuse: the particle is slowing down.
\end{itemize}
Indeed $\vv\cdot\va=v\,a_t$, so its sign is the sign of the tangential acceleration.

\stepn{3} \textbf{$\vomega\cdot\vc{\alpha}$, can be any of the three.} Both are axial, so they are either
parallel or antiparallel; there is no third option for rotation about a fixed axis.
\begin{itemize}
\item \emph{positive} when $\vc{\alpha}\parallel\vomega$: the spin is speeding up;
\item \emph{negative} when antiparallel: the spin is slowing;
\item \emph{zero} when $\vc{\alpha}=\vc0$ (uniform circular motion), or at an instant when
      $\vomega=\vc0$, such as the moment of starting from rest.
\end{itemize}
\ans{$\vomega\cdot\vv=0$ always; the other two may be $+$, $-$ or $0$}
\end{exbox}

\begin{exbox}{When centripetal and tangential accelerations are equal}
A particle starts from rest and moves on a circle with constant angular acceleration
$\alpha=\SI{4}{rad/s^2}$. At what time are its centripetal and tangential accelerations equal in
magnitude?

\stepn{1} \textbf{Write both in terms of $t$.} Starting from rest, $\omega=\alpha t$. Then
\[
a_r=\omega^{2}r=(\alpha t)^{2}r=\alpha^{2}t^{2}r,
\qquad
a_t=\alpha r.
\]

\stepn{2} \textbf{Set them equal.} The radius cancels: a good sign, since the question never gave
one:
\[
\alpha^{2}t^{2}r=\alpha r
\qquad\Longrightarrow\qquad
t^{2}=\frac{1}{\alpha}
\qquad\Longrightarrow\qquad
t=\frac{1}{\sqrt\alpha}=\frac{1}{\sqrt4}=\SI{0.5}{s}.
\]

\stepn{3} \textbf{Interpretation.} Before $t=\SI{0.5}{s}$ the \emph{tangential} part dominates (the
particle is mostly gaining speed); after it, the \emph{centripetal} part takes over and grows without
limit as $t^{2}$. Under constant $\alpha$ the radial part always overtakes the tangential part
eventually. \ans{$t=\SI{0.5}{s}$}
\end{exbox}

\begin{exbox}{Reading the angle between $\vv$ and $\va$}
A particle moves on a circle of radius \SI{54}{m} with speed $v=4t^{2}$ (SI). Find the angle between its
velocity and its total acceleration at $t=\SI{3}{s}$.

\stepn{1} \textbf{Speed:} $v=4(3)^{2}=\SI{36}{m/s}$.

\stepn{2} \textbf{Tangential:} $a_t=\dfrac{dv}{dt}=8t=\SI{24}{m/s^2}$.

\stepn{3} \textbf{Radial:} $a_r=\dfrac{v^{2}}{r}=\dfrac{36^{2}}{54}=\dfrac{1296}{54}=\SI{24}{m/s^2}$.

\stepn{4} \textbf{Angle.} The angle $\phi$ between $\va$ and $\vv$ satisfies
$\tan\phi=\dfrac{a_r}{a_t}=\dfrac{24}{24}=1$, so
\[
\phi=45^\circ.
\]
It is acute, confirming that the particle is speeding up, as the increasing speed $v=4t^{2}$ requires.
\ans{$45^\circ$}
\end{exbox}

\begin{exbox}{Working backwards from the total acceleration}
A particle moves on a circle of radius \SI{2.5}{m}; at one instant its total acceleration
$a=\SI{25}{m/s^2}$ makes $30^\circ$ with the \emph{radius}. Find (a) the radial acceleration, (b) the
speed, (c) the tangential acceleration.

\stepn{1} \textbf{Resolve the given acceleration.} The angle is measured from the radius, so the radial
part takes the cosine and the tangential part the sine:
\[
a_r=a\cos30^\circ=25(0.866)=\SI{21.65}{m/s^2},
\qquad
a_t=a\sin30^\circ=25(0.5)=\SI{12.5}{m/s^2}.
\]

\stepn{2} \textbf{Speed from the radial part.} Since $a_r=v^{2}/r$,
\[
v=\sqrt{a_r r}=\sqrt{21.65\times2.5}=\sqrt{54.1}\approx\SI{7.36}{m/s}.
\]
\ans{$a_r\approx\SI{21.7}{m/s^2}$; $v\approx\SI{7.36}{m/s}$; $a_t=\SI{12.5}{m/s^2}$}

\smallskip
\textbf{Read the angle before you resolve.} Had the $30^\circ$ been measured from the \emph{tangent}
instead, the sine and cosine would swap and every answer would change. Always check which line the
angle is drawn from: this is the commonest error in this type of problem.
\end{exbox}

\begin{exbox}{Average velocity versus instantaneous speed}
A particle moves at constant speed on a circle. Find the ratio of the magnitude of its \emph{average
velocity} to its \emph{instantaneous speed} while it turns through $\theta=\tfrac{\pi}{2}$.

\stepn{1} \textbf{Average velocity needs the displacement}, which is the \emph{chord}, not the arc. For
a quarter circle the chord joining the ends is the hypotenuse of a right isosceles triangle of legs
$r$:
\[
|\Delta\vr|=r\sqrt2.
\]

\stepn{2} \textbf{The time taken} to sweep $\theta$ at angular speed $\omega$ is $t=\theta/\omega$.

\stepn{3} \textbf{Form the ratio.} With instantaneous speed $v=\omega r$,
\[
\frac{|\vv_{\text{avg}}|}{v}
=\frac{r\sqrt2\big/(\theta/\omega)}{\omega r}
=\frac{\sqrt2}{\theta}
=\frac{\sqrt2}{\pi/2}
=\boxed{\frac{2\sqrt2}{\pi}\approx0.90}
\]

\stepn{4} \textbf{Why it must be less than 1.} The chord is always shorter than the arc, so the average
velocity can never match the (constant) speed. Over a \emph{full} circle the displacement is zero and
the ratio falls to $0$, even though the speed never changed: a sharp reminder that average velocity and
average speed are different quantities. \ans{$2\sqrt2/\pi\approx0.90$}
\end{exbox}

\section{Summary}\label{sec:summary6}
Velocity is a vector, so it can change in two ways: in magnitude and in direction. Writing $\vv=v\,\uhat{t}$ and differentiating with the product rule separates the two at once:
$\dfrac{dv}{dt}\,\uhat{t}$ changes the speed, and $v\,\dfrac{d\uhat{t}}{dt}$ changes the direction. Because a unit vector can only rotate,
its rate of change is perpendicular to it, so the two components are automatically perpendicular and
combine by Pythagoras. The table below collects the working relations.

\medskip
{\small\renewcommand{\arraystretch}{1.45}
\begin{center}
\begin{tabular}{@{}lll@{}}
\toprule
Quantity & Formula & Note\\
\midrule
Arc / angle        & $s=r\theta$                                   & $\theta$ in radians, always\\
Linear speed       & $v=r\omega$, \; $\vv=\vomega\ex\vr$           & $\vomega$ along the axis\\
Tangential accel.  & $a_t=\dfrac{dv}{dt}=r\alpha$                  & changes the \emph{speed}; may be $0$\\
Radial accel.      & $a_r=\dfrac{v^{2}}{r}=\omega^{2}r$            & changes the \emph{direction}; $0$ only if $v=0$\\
Total              & $a=\sqrt{a_t^{2}+a_r^{2}}$, \; $\tan\phi=\dfrac{a_r}{a_t}$ & $\phi$ from the tangent\\
Radius of curvature& $R=\dfrac{v^{2}}{a_n}$                        & works on \emph{any} curved path\\
Period / frequency & $T=\dfrac{2\pi}{\omega}=\dfrac1\nu$, \; $\omega=\dfrac{2\pi N}{60}$ & $N$ in rpm\\
Constant $\alpha$  & $\omega=\omega_0+\alpha t$,\; $\Delta\theta=\omega_0t+\tfrac12\alpha t^{2}$ & only if $\alpha$ is constant\\
                   & $\omega^{2}=\omega_0^{2}+2\alpha\,\Delta\theta$ & \\
Varying $\alpha$   & $\displaystyle\int\omega\,d\omega=\int\alpha\,d\theta$ & the $v\,dv/dx$ twin\\
\bottomrule
\end{tabular}
\end{center}
}

\medskip
\textbf{Diagnosing a motion from one angle.} With $\phi$ the angle between $\va$ and $\vv$:
acute $\Rightarrow$ speeding up; $90^\circ$ $\Rightarrow$ uniform; obtuse $\Rightarrow$ slowing.
Equivalently $\vv\cdot\va=v\,a_t$ carries the sign.

\medskip
\textbf{In the plane:} $\vv$, $\va$, $a_t\uhat{t}$, $a_r\uhat{n}$. \quad
\textbf{Along the axis:} $\vomega$, $\vc{\alpha}$. \quad
Hence $\vomega\cdot\vv=0$ always.

\section{Problem Types}\label{sec:types6}
Circular kinematics produces five recurring question types, ranging from a direct calculation of the
centripetal acceleration to the total acceleration of a body that is speeding up as it turns. In every
case, convert any rate given in revolutions per minute to radians per second before doing anything else,
and keep the tangential and radial components separate until the very last step.

\archq{Find the centripetal acceleration}
{``A particle moves in a circle of radius $r$ at constant speed $v$ (or period $T$, or frequency
$\nu$); find its acceleration.''}
{Convert whatever is given into $v$ or $\omega$ first ($v=\omega r$, $\omega=2\pi/T=2\pi\nu$), then
$a_r=\dfrac{v^2}{r}=\omega^2 r$, directed \emph{towards the centre}.}
{Saying the acceleration is zero because the speed is constant. The speed is constant, but the
\emph{velocity} is not, because its direction turns; there is an acceleration, perpendicular to $\vv$.}
{A particle moves on a circle of radius \SI{2}{m} at a constant \SI{4}{m/s}. Its acceleration is:}
{zero}{\SI{2}{m/s^2}}{\SI{16}{m/s^2}}{\SI{8}{m/s^2}\archans{D}{$a_r=v^2/r=16/2=\SI{8}{m/s^2}$, directed
to the centre. Option (A) is the classic error of confusing constant \emph{speed} with constant
velocity.}}

\archq{Convert rpm to angular and linear speed}
{``A wheel or fan rotates at $N$ revolutions per minute; find $\omega$, $v$ or $a_r$.''}
{Convert first: $\omega=\dfrac{2\pi N}{60}$~\si{rad/s}. Then $v=\omega r$ and $a_r=\omega^2r=v^2/r$. One
revolution is $2\pi$ radians, never $360$ in a formula.}
{Leaving $N$ in rpm, or using degrees. Every rotational formula assumes radians per second.}
{A wheel rotates at \SI{60}{rpm}. Its angular speed is:}
{$2\pi$~\si{rad/s}}{$\pi$~\si{rad/s}}{\SI{60}{rad/s}}{$120\pi$~\si{rad/s}\archans{A}{$\omega=\dfrac{2\pi(60)}{60}=2\pi$~\si{rad/s}, that is, one revolution per second.}}

\archq{How $a_r$ scales with $v$, $r$, $T$}
{``If the speed is doubled (or the radius halved), what happens to the centripetal acceleration?''}
{Pick the form that matches what is held fixed: $a_r=\dfrac{v^2}{r}$ when $v$ is given,
$a_r=\omega^2r$ when $\omega$ is given, $a_r=\dfrac{4\pi^2r}{T^2}$ when the period is. These scale
\emph{differently}: at constant $v$, $a_r\propto1/r$; at constant $\omega$, $a_r\propto r$.}
{Using the wrong form, then getting the direction of the dependence on $r$ backwards.}
{Two particles move in circles with the \emph{same speed}, one on a radius twice the other's. The ratio
of their centripetal accelerations, larger circle to smaller, is:}
{$2:1$}{$4:1$}{$1:4$}{$1:2$\archans{D}{At constant $v$, $a_r=v^2/r\propto1/r$. Doubling $r$
\emph{halves} $a_r$, so the ratio is $1:2$.}}

\archq{Angular kinematics with constant $\alpha$}
{``A wheel accelerates from rest to $\omega$ in $t$ seconds; find $\alpha$, the angle turned, or the
number of revolutions.''}
{Use the constant-$\alpha$ equations: $\omega=\omega_0+\alpha t$, $\Delta\theta=\omega_0t+\tfrac12\alpha t^2$,
$\omega^2=\omega_0^2+2\alpha\,\Delta\theta$. To convert an angle to revolutions, divide by $2\pi$.}
{Reporting $\theta$ in radians when revolutions were asked (or vice versa).}
{A wheel starts from rest with $\alpha=\SI{2}{rad/s^2}$. The angle turned in \SI{3}{s} is:}
{\SI{6}{rad}}{\SI{9}{rad}}{\SI{18}{rad}}{\SI{36}{rad}\archans{B}{$\Delta\theta=\tfrac12\alpha t^2
=\tfrac12(2)(9)=\SI{9}{rad}$. Option (A) is $\omega=\alpha t$, an angular \emph{velocity}.}}

\archq{Total acceleration when the speed also changes}
{``A particle speeds up while moving on a circle; find its total acceleration.'' The tell is that the
speed is \emph{not} constant.}
{There are now two perpendicular components: tangential $a_t=r\alpha$ (changing the speed) and
centripetal $a_r=v^2/r$ (changing the direction). The total is
$a=\sqrt{a_t^2+a_r^2}$.}
{Reporting only $a_r$. That is complete only in \emph{uniform} circular motion, where $a_t=0$.}
{A particle on a circle has $a_r=\SI{3}{m/s^2}$ and $a_t=\SI{4}{m/s^2}$. Its total acceleration is:}
{\SI{5}{m/s^2}}{\SI{1}{m/s^2}}{\SI{3.5}{m/s^2}}{\SI{7}{m/s^2}\archans{A}{The components are
perpendicular, so $a=\sqrt{3^2+4^2}=\SI{5}{m/s^2}$.}}

\section{Practice Problems}\label{sec:practice6}

The problems below are of standard examination difficulty and are meant to be worked without looking back at the chapter. Answers, with brief working, follow the list.

\begin{enumerate}[itemsep=3pt]
\item Is the acceleration of a particle in \emph{uniform} circular motion constant or variable? Explain
      in one line.
\item A particle moves on a circle. Which of these may stay constant: (i) velocity,
      (ii) speed, (iii) acceleration, (iv) magnitude of acceleration?
\item A particle moves on a circle of radius \SI{1.0}{cm} with $v=2t$ (cm/s, s). At $t=\SI{1}{s}$ find
      (a) the radial acceleration, (b) the tangential acceleration, (c) the net acceleration.
\item A wheel turns at \SI{300}{rpm}. Find its angular speed in \si{rad/s}.
\item A fan slows uniformly from \SI{900}{rpm} to \SI{300}{rpm} in \SI{5.0}{s}. Find (a) the angular
      acceleration and (b) the number of revolutions made in that time.
\item A particle in circular motion has $a_t=\SI{3}{m/s^2}$ and $a_r=\SI{4}{m/s^2}$. Find the magnitude
      of the net acceleration and the angle it makes with the velocity.
\end{enumerate}

\subsection*{Answers}
{\small
\begin{enumerate}[itemsep=2pt]
\item \emph{Variable}: its magnitude $v^2/r$ is fixed but its direction continuously turns to
stay pointed at the centre, and a vector of changing direction is not constant.
\item (ii) speed and (iv) magnitude of acceleration (both constant in uniform circular motion).
Velocity cannot (its direction turns), and neither can acceleration as a vector, because its radial
part always points at the centre and so turns with the particle. (On a general curved path the
acceleration vector \emph{can} be constant: a projectile is the example.)
\item $v=\SI{2}{cm/s}$, so (a) $a_r=v^2/r=\SI{4.0}{cm/s^2}$; (b) $a_t=dv/dt=\SI{2.0}{cm/s^2}$;
(c) $a=\sqrt{16+4}=2\sqrt5\approx\SI{4.5}{cm/s^2}$.
\item $\omega=\dfrac{2\pi(300)}{60}=10\pi\approx\SI{31.4}{rad/s}$.
\item $\omega_1=30\pi$, $\omega_2=10\pi$~\si{rad/s}. (a) $\alpha=\dfrac{10\pi-30\pi}{5}=-4\pi
\approx\SI{-12.6}{rad/s^2}$. (b) $\theta=\bar\omega t=\dfrac{30\pi+10\pi}{2}(5)=100\pi$~\si{rad}
$=\mathbf{50}$ revolutions.
\item $a=\sqrt{3^2+4^2}=\SI{5}{m/s^2}$, at $\tan^{-1}(4/3)=53^\circ$ to the velocity (acute, so
the particle is speeding up).
\end{enumerate}
}


\section*{Looking Back on Part A}
\addcontentsline{toc}{section}{Looking Back on Part A}

Part A built the language in which the rest of mechanics is written, and it is worth seeing that
language as a single thing rather than a sequence of chapters. Its grammar is the vector, a quantity with
size and direction that adds by the triangle law. Its central technique is resolution: break a vector
into perpendicular components, treat each component as a one-dimensional problem, and recombine. With
that technique the kinematics of a straight line, developed in \chref{ch:straightline}, extends without
change to projectiles in \chref{ch:projectiles} and to relative motion in \chref{ch:relative}.

\begin{itemize}
\item \textbf{Mathematical toolkit (\chref{ch:toolkit}).} The derivative measures a rate, the integral a
      total; dot and cross products, the small-angle approximation and the quadratic formula recur
      throughout.
\item \textbf{Measurement (\chref{ch:measurement}).} Every quantity carries a unit and an uncertainty. Dimensions check a
      formula and can find its shape; relative errors add for products and quotients.
\item \textbf{Straight-line motion (\chref{ch:straightline}).} Velocity and acceleration are rates, slopes on a graph; their
      accumulations are areas. The three kinematic equations hold only for constant acceleration.
\item \textbf{Vectors (\chref{ch:vectors}).} Resolution into components is the master technique of mechanics.
\item \textbf{Projectiles (\chref{ch:projectiles}).} Horizontal and vertical motions are independent and
      share only the clock; the level-ground formulas fail for a launch from a height.
\item \textbf{Relative motion (\chref{ch:relative}).} $\vv_{AB}=\vv_A-\vv_B$, drawn as a triangle;
      minimum time and minimum drift are different goals.
\item \textbf{Circular kinematics (\chref{ch:circkin}).} A changing direction is an acceleration. On any curve the
      acceleration has a tangential part that changes the speed and a radial part, $v^2/r$, that changes the
      direction.
\end{itemize}

\noindent\textbf{The thread into Part B.} Everything so far has described motion without asking what causes
it. \chref{ch:circkin} in particular established that a body moving in a circle \emph{must} have an inward
acceleration. Part B supplies the cause: forces, and the law that connects them to acceleration.
